% =====================================================================
%  Citacoes e referencias: biblatex + biber
%
%  Escolhido no scaffold com --bib=biblatex. Exige os pacotes biblatex e
%  biblatex-abnt e o binario biber. Rode `latexkit doctor` para conferir
%  o que falta na sua instalacao.
%
%  Como citar:
%    \parencite{silva2020}    -> (SILVA, 2020)
%    \textcite{silva2020}     -> Silva (2020), para citacao no corpo da frase
%    \cite{silva2020}
% =====================================================================

\usepackage[
  style=abnt,
  backend=biber,
  backref=false,
  giveninits=true,
  repeatfields=true
]{biblatex}

\addbibresource{bib/references.bib}

% A abntex2cite oferece \citeonline e \apud, e sao esses os nomes que a
% literatura brasileira usa. Reproduzimos os dois sobre os comandos do
% biblatex para que o texto do documento nao precise saber qual backend esta
% em uso: o mesmo content/*.tex compila com abntex2cite e com biblatex.
%
% O par \providecommand + \renewcommand cobre os dois casos: a classe abntex2
% ja define \apud, e redefinir o que existe exige \renewcommand; o beamer nao
% define nenhum dos dois, e ai o \providecommand cria o nome antes.
\providecommand{\citeonline}{}
\renewcommand{\citeonline}[2][]{\textcite[#1]{#2}}
\providecommand{\apud}{}
\renewcommand{\apud}[2]{\parencite{#1} apud \parencite{#2}}

% Imprime a secao de referencias. O main.tex chama esta macro; assim os
% dois backends de bibliografia expoem a mesma interface.
\newcommand{\lkPrintBibliography}{%
  \printbibliography[title=REFER\^ENCIAS]%
}
